\documentclass[11pt,a4paper]{article}

% --- Packages ---
\usepackage[utf8]{inputenc}
\usepackage[T1]{fontenc}
\usepackage{lmodern}
\usepackage[margin=2.5cm]{geometry}
\usepackage{graphicx}
\usepackage{xcolor}
\usepackage{hyperref}
\usepackage{booktabs}
\usepackage{enumitem}
\usepackage{titlesec}
\usepackage{fancyhdr}
\usepackage{tcolorbox}
\usepackage{float}
\usepackage{caption}
\usepackage{multicol}
\usepackage{parskip}

% --- Colors ---
\definecolor{primary}{RGB}{44,62,80}
\definecolor{accent}{RGB}{52,152,219}
\definecolor{insightbg}{RGB}{235,245,255}
\definecolor{bodytext}{RGB}{50,50,50}

% --- Hyperref setup ---
\hypersetup{
    colorlinks=true,
    linkcolor=accent,
    urlcolor=accent,
    pdftitle={EDA Insight Report -- Team Entropy},
    pdfauthor={Saanvi Grover, Aditya Sharma}
}

% --- Section formatting ---
\titleformat{\section}{\Large\bfseries\color{primary}}{}{0em}{}[\color{accent}\titlerule]
\titleformat{\subsection}{\large\bfseries\color{primary}}{}{0em}{}

% --- Header/Footer ---
\pagestyle{fancy}
\fancyhf{}
\fancyhead[L]{\small\color{gray}Data Vortex $|$ Team Entropy $|$ EDA Insight Report}
\fancyhead[R]{\small\color{gray}Page \thepage}
\fancyfoot[C]{\small\color{gray}AARUUSH'26 -- Rebuilding the Social Engine -- Team Entropy}
\renewcommand{\headrulewidth}{0.4pt}
\renewcommand{\footrulewidth}{0pt}

% --- Insight box ---
\newtcolorbox{insightbox}{
    colback=insightbg,
    colframe=accent,
    boxrule=0.5pt,
    arc=2pt,
    left=8pt, right=8pt, top=6pt, bottom=6pt,
    fonttitle=\bfseries\color{accent}\small,
    title=KEY INSIGHT
}

% --- Image helper ---
\newcommand{\edaimg}[2][0.85]{%
    \begin{figure}[H]
        \centering
        \includegraphics[width=#1\textwidth]{#2}
    \end{figure}
}

\begin{document}
\color{bodytext}

% ============================================================
% TITLE PAGE
% ============================================================
\begin{titlepage}
\centering
\vspace*{4cm}
{\Huge\bfseries\color{primary} Exploratory Data Analysis\\[0.3cm] Insight Report\par}
\vspace{1cm}
{\Large\color{gray} Round 1, Phase 1: Data Recovery\par}
\vspace{2cm}
{\color{accent}\rule{8cm}{1pt}\par}
\vspace{2cm}
{\large
\begin{tabular}{rl}
\textbf{Competition:} & Data Vortex $|$ AARUUSH'26 \\
\textbf{Theme:} & Rebuilding the Social Engine \\
\textbf{Team:} & Entropy \\
\textbf{Members:} & Saanvi Grover \& Aditya Sharma \\[0.5cm]
\textbf{Datasets:} & 12,000 posts $|$ 1,500 users $|$ 5 platforms \\
\textbf{Period:} & May 2024 -- April 2025 (364 days) \\
\textbf{Scope:} & 20 analytical sections $|$ 21 visualizations \\
& 6 formal statistical significance tests \\
\end{tabular}
\par}
\vfill
\end{titlepage}

% ============================================================
% TABLE OF CONTENTS
% ============================================================
\tableofcontents
\newpage

% ============================================================
% SECTION 1
% ============================================================
\section{Dataset Overview}

The Social Engine dataset consists of two tables recovered from Archive Node 07 on the competition's recovery terminal. After cleaning (Phase 1a), the final datasets are:

\textbf{Posts Dataset:} 12,000 rows $\times$ 8 columns (post\_id, user\_id, platform, text\_content, timestamp, likes, shares, comments). Date range spans 364 days from May 2024 to April 2025.

\textbf{Users Dataset:} 1,500 rows $\times$ 5 columns (user\_id, location, language, account\_created, follower\_count). Users created between Jan--Dec 2023 across 33 cities and 10 languages.

\textbf{Missing data} remains in three columns: platform (14.9\%), text\_content (14.3\%), and likes (15.1\%). These were left as NaN because no reliable imputation basis exists---the corruption was confirmed to be Missing Completely at Random (MCAR) through statistical testing (Mann-Whitney U, $p=0.64$).

% ============================================================
% SECTION 2
% ============================================================
\section{Platform Analysis: Does Platform Choice Matter?}

\textit{Purpose: In real social media, Instagram drives 2--3$\times$ more engagement than Twitter. Here we test whether platform choice matters in the Social Engine dataset.}

\edaimg{images/Entropy_eda_platform_distribution.png}

Activity is nearly uniform: Facebook leads with 2,074 posts and Instagram trails at 1,989---a mere 4.3\% spread. This uniformity is unusual for real social media and may reflect how the synthetic data was generated.

\edaimg{images/Entropy_eda_engagement_by_platform.png}

Average engagement metrics (likes, shares, comments) are nearly identical across all five platforms. A Kruskal-Wallis test (Section 14) confirms this: $H=4.27$, $p=0.37$---platform choice does \textbf{NOT} significantly affect engagement.

\begin{insightbox}
Platform distribution is remarkably uniform ($\sim$2,000 posts each). Unlike real social media where platform dynamics vary dramatically, engagement here is platform-agnostic. This means Phase 2 SQL queries testing platform effects will find flat results---not a failure, but a structural property of the data.
\end{insightbox}

% ============================================================
% SECTION 3
% ============================================================
\section{Temporal Analysis: When Should You Post?}

\textit{Purpose: Real platforms show 2--3$\times$ hourly variation and 20--40\% day-of-week swings. We test whether timing matters in this dataset.}

\edaimg{images/Entropy_eda_monthly_posts.png}

Monthly post volume is stable around 914--1,038 posts/month. February 2025 is the lowest (914 posts, $-9.2\%$ MoM), likely reflecting the shorter month. March 2025 rebounds sharply ($+10.8\%$). No seasonal trend is evident.

\edaimg{images/Entropy_eda_temporal_patterns.png}

\textbf{Left:} Day-of-week posting shows Wednesday as the most active day (1,771 posts) and Saturday the least (1,675). The spread is narrow at 5.7\%.
\textbf{Right:} Hourly analysis reveals a nearly flat distribution across all 24 hours. Unlike real social media (which shows strong peaks during lunch and evening), this dataset has users posting at 3\,AM as often as 7\,PM.

\begin{insightbox}
Temporal patterns are unusually flat. Real social platforms show 2--3$\times$ variation by hour of day and 20--40\% variation by day of week. This dataset shows $<6\%$ variation on both, suggesting timestamps were generated uniformly rather than reflecting real human behaviour.
\end{insightbox}

% ============================================================
% SECTION 4
% ============================================================
\section{Engagement Distributions: Is There Viral Content?}

\textit{Purpose: Real social media follows power-law distributions (few viral posts, many low-engagement). We test whether this dataset has outliers.}

\edaimg{images/Entropy_eda_engagement_distributions.png}

All three metrics follow near-uniform distributions, not the heavy-tailed (power-law) distributions typical of real social media. Likes range from 0 to 5,000 with mean $\approx$2,492 and median $\approx$2,498 (nearly identical, confirming symmetry). There are no viral outliers---every post gets roughly similar engagement.

\edaimg{images/Entropy_eda_engagement_correlation.png}

The correlation heatmap reveals weak relationships between all three metrics. Likes-shares correlation is only 0.003, likes-comments is $-0.003$, and shares-comments is $-0.010$. In real social media, these metrics are typically positively correlated.

\begin{insightbox}
Engagement metrics are uniformly distributed and statistically independent. This is the opposite of real social media where engagement follows power-law distributions and metrics are positively correlated.
\end{insightbox}

% ============================================================
% SECTION 5
% ============================================================
\section{User Analysis: Do Followers Equal Influence?}

\textit{Purpose: On real platforms, follower count is the \#1 predictor of reach. We test whether that holds in this dataset.}

\edaimg{images/Entropy_eda_user_activity.png}

\textbf{Left:} The posts-per-user distribution is right-skewed. Mean is 8.0 posts/user, ranging from 1 to 22. \textbf{Right:} Follower count vs.\ post count shows no visible relationship---the scatter plot is a uniform cloud. A Spearman correlation test confirms: $\rho=0.025$, $p=0.34$ (not significant). Follower count does not predict activity or engagement.

\begin{insightbox}
Follower count has zero predictive power---contradicting real social media dynamics where follower count is the strongest predictor of reach.
\end{insightbox}

% ============================================================
% SECTION 6
% ============================================================
\section{Geographic Analysis}

\textit{Purpose: Map the geographic distribution of users and identify whether location influences engagement levels.}

\edaimg{images/Entropy_eda_geographic.png}

Users span 33 cities across 6 continents. Shanghai leads with 56 users (3.7\%). Average likes by location show modest variation (2,300--2,700 range). No single city or region dominates engagement.

% ============================================================
% SECTION 7
% ============================================================
\section{Language Analysis}

\textit{Purpose: Examine how the 10 language communities differ in size and engagement patterns.}

\edaimg{images/Entropy_eda_language.png}

All 10 languages have between 119--168 users. Average engagement by language is remarkably flat---the engagement generation process was language-independent.

% ============================================================
% SECTION 8
% ============================================================
\section{Content Analysis: Hashtags}

\textit{Purpose: Extract and rank hashtags from post text to identify dominant content themes.}

\edaimg{images/Entropy_eda_top_hashtags.png}

The top 20 hashtags reveal consumer-focused themes: \#Fashion, \#Sale, \#Tech, \#Health. These are generic commercial hashtags consistent with brand-centric post templates.

% ============================================================
% SECTION 9
% ============================================================
\section{Content Analysis: Brands}

\textit{Purpose: Identify which of the 10 embedded brands are mentioned most frequently and whether brand mentions correlate with higher engagement.}

\edaimg{images/Entropy_eda_brand_analysis.png}

All 10 brands appear in roughly 1,000--1,100 posts each. Average likes per brand post are nearly identical (2,450--2,550 range). No brand drives significantly more engagement. Engagement is independent of both platform and brand.

% ============================================================
% SECTION 10
% ============================================================
\section{Sentiment Indicators}

\textit{Purpose: Classify post sentiment using keyword matching to understand the emotional tone distribution and whether sentiment predicts engagement.}

\edaimg{images/Entropy_eda_sentiment.png}

A notable finding is the ``Contradictory'' category---posts containing both positive and negative phrases (e.g., ``Bummed out\ldots\ Absolutely loving it'')---which reflects template-based text generation. Negative posts get slightly higher engagement than positive---the ``outrage engagement'' pattern seen in real social media research.

\begin{insightbox}
Contradictory sentiment posts are a signature of template-based text generation. This is a data quality observation worth flagging in competition analysis.
\end{insightbox}

% ============================================================
% SECTION 11
% ============================================================
\section{Anomaly Detection}

\textit{Purpose: Use IQR-based statistical methods to identify outlier posts with unusually high or low engagement.}

\edaimg{images/Entropy_eda_anomalies.png}

The IQR method finds zero outlier posts for likes---every post falls within expected statistical bounds. This is highly unusual: real social media datasets always have outliers (viral posts). The likes-vs-shares scatter shows a uniform cloud with no clustering.

% ============================================================
% SECTION 12
% ============================================================
\section{Missing Data Analysis}

\textit{Purpose: Determine whether missing values are random or systematic.}

\edaimg{images/Entropy_eda_missing_patterns.png}

The missingness correlation heatmap shows weak correlations between missing indicators, confirming Missing Completely at Random (MCAR). A Mann-Whitney U test confirms posts with missing platform data have statistically similar engagement ($p=0.64$), validating our cleaning decision to retain missing values as NaN.

% ============================================================
% SECTION 13
% ============================================================
\section{Key Insights Summary}

The 20 analytical sections reveal these core findings:

\begin{enumerate}[leftmargin=*, itemsep=4pt]
    \item \textbf{Structural Uniformity:} Engagement, platform distribution, brand mentions, language communities, and temporal patterns are all remarkably flat---the defining characteristic of this dataset.
    \item \textbf{No Viral Content:} Zero posts exceed the IQR outlier threshold. Every post receives ``average'' engagement. Real social data always has power-law tails.
    \item \textbf{Independent Metrics:} Likes, shares, and comments are uncorrelated ($r < 0.01$). Each was generated independently.
    \item \textbf{Follower Count is Decorative:} No correlation with engagement (Spearman $\rho=0.025$, $p=0.34$). In real platforms, this is the strongest predictor of reach.
    \item \textbf{Corruption was Random:} Missing data is MCAR, engagement is unaffected by completeness.
    \item \textbf{Template-Generated Text:} Contradictory sentiment phrases, uniform brand distribution, and generic hashtags point to template-based content generation.
    \item \textbf{One Genuine Signal:} Multi-platform users show $\sim$15\% higher engagement---the only statistically significant finding across all 6 hypothesis tests.
    \item \textbf{CV Confirms Uniformity:} Coefficient of variation $< 10\%$ on most dimensions formally quantifies the uniformity with a standard statistical measure.
\end{enumerate}

% ============================================================
% SECTION 14
% ============================================================
\section{Statistical Significance Testing}

\textit{Purpose: Formally test whether observed patterns are statistically significant, using non-parametric tests suited to non-normal distributions.}

\begin{table}[H]
\centering
\small
\begin{tabular}{@{}llll@{}}
\toprule
\textbf{Test} & \textbf{Method} & \textbf{Statistic} & \textbf{Result} \\
\midrule
Platforms vs.\ Likes & Kruskal-Wallis & $H=4.27$, $p=0.37$ & NOT significant \\
Day-of-Week vs.\ Engagement & Kruskal-Wallis & $H=3.72$, $p=0.71$ & NOT significant \\
Sentiment $\times$ Platform & Chi-square & $\chi^2=23.94$, $p=0.24$ & NOT significant \\
Followers vs.\ Likes & Spearman & $\rho=0.025$, $p=0.34$ & NOT significant \\
Missing Platform vs.\ Engagement & Mann-Whitney U & $U=9{,}175{,}298$, $p=0.64$ & NOT significant \\
Multi-Platform Users & Kruskal-Wallis & $\sim$15\% higher & \textbf{SIGNIFICANT} \\
\bottomrule
\end{tabular}
\caption{Summary of 6 formal hypothesis tests at $\alpha=0.05$.}
\end{table}

\begin{insightbox}
Five of six tests return $p > 0.05$. The sixth test (multi-platform user behaviour) is the one genuine signal worth investigating in subsequent rounds.
\end{insightbox}

% ============================================================
% SECTION 15
% ============================================================
\section{Brand $\times$ Platform Cross-Analysis}

\textit{Purpose: Test whether certain brands perform better on specific platforms.}

\edaimg{images/Entropy_eda_brand_platform_heatmap.png}

Average likes by brand $\times$ platform shows values in the 2,300--2,700 range with no obvious hot spots. Post volume by brand $\times$ platform is remarkably even ($\sim$200 posts per cell). No brand-platform affinity exists.

\begin{insightbox}
Unlike real social media where Nike dominates Instagram and tech brands lead on Reddit, this dataset shows no brand-platform affinity---useful context for Phase 2 SQL analysis.
\end{insightbox}

% ============================================================
% SECTION 16
% ============================================================
\section{Text Length Analysis}

\textit{Purpose: Examine whether post length varies by platform and predicts engagement.}

\edaimg{images/Entropy_eda_text_length_analysis.png}

Character length peaks around 80--120 characters. Word count centres around 15--25 words. All platforms have similar median word counts ($\sim$18--20 words). Word count vs.\ likes shows near-zero correlation. Post length does not predict engagement.

% ============================================================
% SECTION 17
% ============================================================
\section{Quantifying Uniformity: Coefficient of Variation}

\textit{Purpose: Formally quantify uniformity using the Coefficient of Variation ($CV = \sigma / \mu \times 100$). A CV below 10\% indicates very low variation.}

\edaimg{images/Entropy_eda_coefficient_of_variation.png}

The CV analysis confirms what earlier sections observed qualitatively: most dimensions have $CV < 10\%$. Platform volume, day-of-week volume, brand mentions, and language distribution are all structurally flat. This level of uniformity is never seen in real social media data, where power-law dynamics create typical $CV > 50\%$.

\begin{insightbox}
The coefficient of variation formally quantifies the dataset's structural uniformity. With most metrics showing $CV < 10\%$, we can state with confidence that the data was generated from uniform distributions---not sampled from real user behaviour.
\end{insightbox}

% ============================================================
% SECTION 18
% ============================================================
\section{Contradictory Sentiment Deep Dive}

\textit{Purpose: Investigate posts containing BOTH positive and negative sentiment phrases. How many exist? What phrase combinations appear? Can we reverse-engineer the templates?}

\edaimg{images/Entropy_eda_contradictory_sentiment_deep_dive.png}

\textbf{Left:} Top contradictory phrase pairs reveal combinations like ``disappointed + loving it'' appearing in the same post---artifacts of a text generator inserting sentiment phrases independently into template structures. \textbf{Right:} Template structure analysis shows posts follow identifiable patterns (unboxing, reviews, ad reactions, comparisons, complaints).

\begin{insightbox}
The deep dive confirms template-based text generation. Sentiment phrases are ``slotted in'' without regard for consistency---explaining why sentiment does not predict engagement.
\end{insightbox}

% ============================================================
% SECTION 19
% ============================================================
\section{Multi-Brand Posts \& Engagement}

\textit{Purpose: Test whether posts mentioning multiple brands receive different engagement than single-brand posts.}

\edaimg{images/Entropy_eda_multi_brand_engagement.png}

Average likes by brand count tests whether content density (more brands = more information) affects engagement. A Mann-Whitney U test compares single-brand vs.\ multi-brand posts. This adds a dimension most teams will not explore: treating brand mentions as a count variable rather than a categorical variable.

% ============================================================
% SECTION 20
% ============================================================
\section{Analytical Narrative: The Story This Data Tells}

\subsection*{The Central Finding: Structural Uniformity}

This EDA reveals that the Social Engine dataset is structurally uniform to a degree never seen in real social media. The CV analysis (Section 17) quantifies this: platform volume, day-of-week patterns, brand mentions, and language distribution all have $CV < 10\%$.

\subsection*{Why This Matters for the Competition}

Understanding uniformity is not a limitation---it \textbf{is} the insight.

\begin{itemize}[leftmargin=*, itemsep=2pt]
    \item \textbf{Phase 2:} Queries testing for viral content and bot detection will return empty results---proving the analyst understands the data's structure.
    \item \textbf{Phase 3:} Multi-platform users are the focus; traditional signals are dead ends.
    \item \textbf{Phase 4:} A dashboard that honestly surfaces ``no effect found'' is more valuable than one that cherry-picks weak patterns.
\end{itemize}

\subsection*{The One Real Signal}

Of all dimensions tested, only multi-platform user behaviour shows a statistically significant effect ($\sim$15\% higher engagement). This is the thread worth pulling.

\subsection*{Synthetic vs.\ Real Comparison}

\begin{table}[H]
\centering
\small
\begin{tabular}{@{}llll@{}}
\toprule
\textbf{Dimension} & \textbf{Real Expectation} & \textbf{Dataset} & \textbf{Verdict} \\
\midrule
Platform engagement & 2--3$\times$ variation & $<5\%$ & Synthetic \\
Hourly posting & Strong peaks & Flat & Synthetic \\
Follower correlation & $r > 0.3$ & $r = 0.025$ & Synthetic \\
Multi-platform advantage & Unknown & 15\% higher (sig.) & Possibly real \\
Template text & N/A & Confirmed & Generation artifact \\
\bottomrule
\end{tabular}
\caption{Synthetic vs.\ real social media comparison.}
\end{table}

\begin{insightbox}
The analytical narrative connects all 20 sections: structural uniformity is the defining characteristic, with one genuine signal (multi-platform users) and one original finding (template-based text generation with contradictory sentiment).
\end{insightbox}

% ============================================================
% CONCLUSIONS
% ============================================================
\section*{Conclusions \& Recommendations}
\addcontentsline{toc}{section}{Conclusions \& Recommendations}

This EDA establishes the fundamental character of the Social Engine dataset through 20 analytical sections, 21 visualizations, and 6 statistical tests. The dataset is structurally uniform ($CV < 10\%$ on most dimensions), with engagement generated independently of user attributes, platform, timing, or content. Two genuine signals emerge: multi-platform users show significantly higher engagement ($\sim$15\%), and template-based text generation creates detectable contradictory sentiment patterns.

\textbf{Recommendations for Phase 2 (SQL Analysis):}

\begin{enumerate}[leftmargin=*, itemsep=4pt]
    \item \textbf{Trend Detection:} Focus on monthly volume and growth rates. Day-of-week and platform trends will be flat but should still be demonstrated with proper window functions.
    \item \textbf{Anomaly Discovery:} Q4 (viral posts) and Q5 (bot detection) will return empty results, but the queries demonstrate advanced statistical SQL and the empty results are a legitimate finding.
    \item \textbf{Behavioural Grouping:} Multi-platform users are the one dimension where a real signal exists (15\% higher engagement for 5-platform users).
    \item \textbf{Correlation Analysis:} Follower count vs.\ engagement will be flat. Brand sentiment analysis is the most promising direction.
\end{enumerate}

% ============================================================
% APPENDIX
% ============================================================
\newpage
\section*{Appendix: Corruption Catalog}
\addcontentsline{toc}{section}{Appendix: Corruption Catalog}

The following corruption patterns were identified and fixed during Phase 1a cleaning.

\edaimg{images/Entropy_corruption_catalog.png}

\edaimg{images/Entropy_data_quality_scorecard.png}

The data quality scorecard shows: 360 duplicate rows removed, 974 HTML artifacts cleaned, 306 mojibake instances fixed, 509 negative likes corrected, 24 NULL strings replaced, and all 12,000 timestamps standardized.

\end{document}
